\documentclass{sbc2023}
\usepackage{float}
\usepackage{graphicx}
%\usepackage[utf8]{inputenc} % Geralmente não necessário com fontspec, que assume UTF-8
\usepackage[misc,geometry]{ifsym}
\usepackage{fontspec} % Para compilar com XeLaTeX ou LuaLaTeX
\usepackage{fontawesome}
\usepackage{academicons}
\usepackage{color}
\usepackage{hyperref}
\usepackage{aas_macros}
\usepackage[bottom]{footmisc}
\usepackage{supertabular}
\usepackage{afterpage}
\usepackage{url}
\usepackage{pifont}
\usepackage{multicol}
\usepackage{multirow}
\usepackage{subcaption} 
\usepackage{tabularx}
\usepackage{booktabs}

% REMOVIDO O CONFLITO: biblatex vs natbib
% O sbc2023.cls já carrega natbib, então usamos ele
% \bibliography{references}

% REMOVIDAS as configurações conflitantes do biblatex:
% \usepackage[backend=biber, style=numeric, sorting=none]{biblatex}
% \addbibresource{refs.bib}
% \DeclareCiteCommand{\cite}...

% Cores personalizadas
\definecolor{orcidlogo}{rgb}{0.37,0.48,0.13}
\definecolor{unilogo}{rgb}{0.16, 0.26, 0.58}
\definecolor{maillogo}{rgb}{0.58, 0.16, 0.26}
\definecolor{darkblue}{rgb}{0.0,0.0,0.0}

% Configuração de hiperlinks
\hypersetup{colorlinks,breaklinks,
            linkcolor=darkblue,urlcolor=darkblue,
            anchorcolor=darkblue,citecolor=darkblue}
%\hypersetup{draft} % Manter comentado para ativar os hiperlinks no PDF final

% Metadados do Journal (verifique se SBCS aceita todos estes)
\jid{JBCS}
\jtitle{Journal of the Brazilian Computer Society, 2025, XX:1, }
\doi{10.5753/jbcs.2025.XXXXXX}
\copyrightstatement{This work is licensed under a Creative Commons Attribution 4.0 International License}
\jyear{2025}

\title[High-Dimensional Structures for Image Similarity]{Evaluating High-Dimensional Data Structures for RGB Image Similarity Search: LSH and M-Tree vs. Classical Approaches}

% Autores e Afiliações
\author[Carrieiros, Rocha, Ribeiro, Temponi, Moraes et al. 2025]{
\affil{\textbf{Luan Barbosa Rosa Carrieiros}~\href{https://orcid.org/0009-0007-2310-1129}{\textcolor{orcidlogo}{\aiOrcid}}~\textcolor{blue}{\faEnvelopeO}~~[~\textbf{Pontifical Catholic University of Minas Gerais}~|\href{mailto:luan.rosa@sga.pucminas.br}{~\textbf{\textit{luan.rosa@sga.pucminas.br}}}~]}

\affil{\textbf{Diego Moreira Rocha}~\href{https://orcid.org/0000-0002-7110-2026}{\textcolor{orcidlogo}{\aiOrcid}}~~[~\textbf{Pontifical Catholic University of Minas Gerais}~|\href{mailto:diego.moreira@sga.pucminas.br}{~\textbf{\textit{diego.moreira@sga.pucminas.br}}}~]}

\affil{\textbf{Iago Fereguetti Ribeiro}~\href{https://orcid.org/0000-0003-3052-3016}{\textcolor{orcidlogo}{\aiOrcid}}~~[~\textbf{Pontifical Catholic University of Minas Gerais}~|\href{mailto:iago.fereguetti@sga.pucminas.br}{~\textbf{\textit{iago.fereguetti@sga.pucminas.br}}}~]}

\affil{\textbf{Bernardo Ferreira Temponi}~\href{https://orcid.org/0000-0001-4892-5537}{\textcolor{orcidlogo}{\aiOrcid}}~~[~\textbf{Pontifical Catholic University of Minas Gerais}~|\href{mailto:bernardo.temponi@sga.pucminas.br}{~\textbf{\textit{bernardo.temponi@sga.pucminas.br}}}~]}

\affil{\textbf{Arthur Gonçalves de Moraes}~\href{https://orcid.org/0000-0001-3195-1605}{\textcolor{orcidlogo}{\aiOrcid}}~~[~\textbf{Pontifical Catholic University of Minas Gerais}~|\href{mailto:arthur.moraes@sga.pucminas.br}{~\textbf{\textit{arthur.moraes@sga.pucminas.br}}}~]}
}

\begin{document}

\begin{frontmatter}
\maketitle

\begin{mail}
PUC Minas, Instituto de Ciências Exatas e Informática (ICEI), Av. Dom José Gaspar, 500, Coração Eucarístico, Belo Horizonte, MG, 30535-901, Brazil.
\end{mail}

\begin{abstract}
\textbf{Abstract.} \\
High-dimensional structures (LSH, M-Tree) promise improved similarity search performance. We evaluate seven structures—five classical (Linear, Hash, Hash Dynamic, Octree, Quadtree) plus LSH and M-Tree—on Animals-10 (26,179 images, OpenCV extraction, threshold=40.0). \textbf{Results contradict expectations:} LSH insertion (53.91ms) is 53.9× slower than Linear (1.00ms); classical structures dominate (Octree: 2.12ms search, 100\% precision). M-Tree achieves only 91.3\% recall. This reveals \textit{bidirectional curse of dimensionality}—high-dim structures impose overhead in low-dim RGB 3D space, challenging assumptions about modern algorithm superiority over classical approaches.
\end{abstract}

\begin{keywords}
Image similarity; LSH; M-Tree; Spatial indexing; Curse of dimensionality; RGB; Performance analysis; OpenCV; Animals-10.
\end{keywords}

\end{frontmatter}


\section{Introduction}
\label{sec:intro}

Image similarity search in RGB color space represents a fundamental problem in computer vision with applications spanning content-based image retrieval, visual search engines, and recommendation systems. While RGB space contains only three dimensions, the selection of efficient indexing structures remains an active research area as datasets scale to millions of images.

Classical spatial indexing approaches—including hash tables, octrees, and quadtrees—have demonstrated consistent effectiveness in low-dimensional spaces through geometric partitioning and spatial locality exploitation. However, the emergence of sophisticated high-dimensional indexing structures raises important questions about their applicability beyond their original design context. Locality-Sensitive Hashing (LSH)~\cite{indyk1998approximate} promises sublinear query time through probabilistic approximation using random projections, specifically designed to combat the curse of dimensionality in high-dimensional spaces (d≥50). M-Tree~\cite{ciaccia1997m} leverages metric space properties and triangle inequality pruning to accelerate similarity search while theoretically guaranteeing exact results.

\textbf{Research Question:} These structures were designed for high-dimensional scenarios such as 128-dimensional SIFT descriptors or 512-dimensional deep learning embeddings. Do they provide benefits when applied to low-dimensional RGB 3D space, or does their complexity impose unnecessary overhead?

This study addresses this gap through comprehensive empirical evaluation. We implement and benchmark seven data structures: five classical approaches (Linear Search, Hash Search, Hash Dynamic Search, Octree, Quadtree) and two modern high-dimensional structures (LSH with random projections, M-Tree with triangle inequality pruning). Experiments are conducted on the Animals-10 dataset~\cite{kaggle_animals10} containing 26,179 real animal photographs with RGB extraction via OpenCV~\cite{opencv_library}.

Our findings challenge conventional assumptions: high-dimensional structures impose significant performance and precision penalties in RGB 3D space. LSH insertion (53.91ms) is 53.9× slower than Linear Search (1.00ms), while classical approaches dominate both insertion and search. This reveals that the curse of dimensionality operates \textit{bidirectionally}—specialized high-dimensional structures suffer unnecessary overhead in low-dimensional contexts, contradicting the assumption that algorithmic sophistication universally translates to practical superiority.

\section{Methodology}
\label{sec:methodology}

RGB similarity search retrieves images $I_i=(r_i,g_i,b_i)$ within Euclidean distance threshold $\tau=40.0$ from query $q=(r_q,g_q,b_q)$: $d(q,I_i)=\sqrt{(r_q-r_i)^2+(g_q-g_i)^2+(b_q-b_i)^2} \leq \tau$.

\subsection{Algorithm Descriptions}

\textbf{Classical Structures (5):} Linear Search serves as the baseline with sequential $O(n)$ examination guaranteeing exact results. Hash Search employs 3D spatial grid partitioning with dynamic cell sizing (cell=30) for $O(1)$ expected insertion and $O(k \cdot \rho)$ search where $k$ represents neighboring cells. Hash Dynamic extends this with adaptive concentric expansion from the query cell. Octree implements recursive 3D spatial tree partitioning (maxPerNode=10) with geometric pruning, achieving $O(\log n)$ expected operations. Quadtree uses 2D iterative spatial tree (maxPerNode=25) projecting onto R,G dimensions with blue channel verified during final distance computation.

\textbf{LSH (Locality-Sensitive Hashing)~\cite{indyk1998approximate}:} Implements probabilistic approximate nearest neighbor search through random projections. Uses L=10 independent hash tables, each with k=4 hash functions and projection width w=50.0. Each hash function projects the d-dimensional input onto random directions, partitioning space into buckets. Query processing searches all L tables in parallel, retrieving candidates from matching buckets. Designed for high-dimensional spaces (d≥50) to achieve sublinear $O(n^\rho)$ query time where $\rho < 1$, trading precision for speed through probabilistic guarantees.

\textbf{M-Tree (Metric Tree)~\cite{ciaccia1997m}:} Implements distance-based indexing for metric spaces using triangle inequality pruning. Each node stores entries with covering radius, enabling pruning of entire subtrees during search when triangle inequality proves they cannot contain results. Node capacity=10 triggers balanced splitting. Theoretically guarantees exact search results while reducing distance computations through geometric pruning in metric space.

\subsection{Experimental Setup}

All seven structures implemented in C++17 with unified interface enabling direct performance comparison. The CREATE→TEST→DESTROY memory management pattern ensures fair evaluation by reconstructing each structure from scratch per test, eliminating caching artifacts and reducing memory consumption by over 70\%.

Dataset comprises 26,179 real animal photographs from the Animals-10 Kaggle collection~\cite{kaggle_animals10}, with RGB values extracted via OpenCV~\cite{opencv_library} file-based processing. Approximately 5\% of images produced libpng iCCP warnings regarding incorrect sRGB profiles, but RGB extraction succeeded for all images. Query image RGB(102,101,104) extracted from ./query/query.jpg with similarity threshold $\tau = 40.0$ based on Euclidean distance.

Benchmark evaluates five dataset scales: 1K, 5K, 10K, 25K, and full 26,179 images, measuring insertion time (structure construction), search time (threshold-based query), and precision (number of similar images found).


\section{Results and Analysis}
\label{sec:results}

\subsection{Performance Analysis}

Table~\ref{tab:real_performance_ms} presents comprehensive performance results for all seven data structures across five dataset scales from 1K to 26,179 real images. The results reveal stark performance divergence between classical spatial indexing methods and modern high-dimensional structures when applied to RGB 3D space, challenging conventional assumptions about algorithmic sophistication translating to practical performance gains.

\begin{table}[H]
    \footnotesize
    \centering
    \caption{Performance Comparison on Real Image Datasets (times in milliseconds)}
    \label{tab:real_performance_ms}
    \setlength{\tabcolsep}{2pt}
    \begin{tabular}{ll rrrrr}
        \toprule
        \textbf{Algorithm} & \textbf{Op.} & \textbf{1K} & \textbf{5K} & \textbf{10K} & \textbf{25K} & \textbf{26K} \\
        \midrule
        \multirow{2}{*}{Linear Search} & Insert & \textbf{0.01} & \textbf{0.15} & \textbf{0.29} & \textbf{0.78} & \textbf{1.00} \\
        & Search & 0.36 & 0.36 & 0.77 & 2.19 & 2.32 \\
        \cmidrule(lr){2-7}
        \multirow{2}{*}{Hash Search} & Insert & 0.14 & 0.73 & 1.50 & 3.65 & 3.68 \\
        & Search & \textbf{0.08} & 0.49 & 0.93 & 2.62 & 2.83 \\
        \cmidrule(lr){2-7}
        \multirow{2}{*}{Hash Dynamic} & Insert & 0.14 & 0.83 & 1.58 & 4.07 & 4.14 \\
        & Search & \textbf{0.06} & \textbf{0.34} & 0.77 & 2.02 & 2.15 \\
        \cmidrule(lr){2-7}
        \multirow{2}{*}{Quadtree} & Insert & 0.25 & 0.99 & 2.16 & 6.00 & 6.92 \\
        & Search & \textbf{0.06} & 0.35 & \textbf{0.76} & \textbf{2.00} & 2.25 \\
        \cmidrule(lr){2-7}
        \multirow{2}{*}{Octree} & Insert & 0.15 & 0.87 & 1.89 & 5.35 & 5.48 \\
        & Search & \textbf{0.06} & 0.38 & 1.07 & 2.23 & \textbf{2.12} \\
        \cmidrule(lr){2-7}
        \multirow{2}{*}{LSH} & Insert & 2.34 & 10.23 & 20.29 & 52.20 & 53.91 \\
        & Search & 0.18 & 0.77 & 1.55 & 4.70 & 4.62 \\
        \cmidrule(lr){2-7}
        \multirow{2}{*}{M-Tree} & Insert & 0.30 & 1.62 & 4.26 & 11.22 & 11.85 \\
        & Search & \textbf{0.06} & 0.35 & 0.79 & 2.22 & 2.46 \\
        \bottomrule
    \end{tabular}
    \\[0.1cm]
    \footnotesize\textit{Bold values indicate best performance in each category (Insert/Search) per scale. 26K represents the full Animals-10 dataset (26,179 images).}
\end{table}
\vspace{-0.3cm}

\textbf{Key Performance Findings:}

\textbf{Insertion Performance:} Linear Search achieves the fastest insertion across all scales (1.00ms for 26K images) through simple array append operations. Classical tree structures show moderate overhead (Octree: 5.48ms, Quadtree: 6.92ms). High-dimensional structures exhibit severe degradation: LSH insertion (53.91ms) is 53.9× slower than Linear and 7.8× slower than Quadtree, while M-Tree (11.85ms) is 11.9× slower than Linear. The LSH overhead stems from computing random projections and maintaining L=10 hash tables, revealing impractical construction costs for RGB applications.

\textbf{Search Performance:} Classical structures dominate with Octree achieving 2.12ms, Quadtree 2.25ms, and Hash Dynamic 2.15ms—all within 6\% performance band. Linear Search (2.32ms) remains competitive despite $O(n)$ complexity, benefiting from cache-friendly sequential access. LSH search (4.62ms) is 2.18× slower than Octree despite theoretical sublinear query promises, invalidating its core value proposition. M-Tree (2.46ms) provides no advantage over simpler Octree while sacrificing precision.

\textbf{Bidirectional Curse of Dimensionality:} Results demonstrate that high-dimensional structures suffer not only in truly high-dimensional spaces (their original design motivation) but also impose unnecessary overhead in low-dimensional RGB 3D contexts. The curse operates bidirectionally—structures optimized for d≥50 dimensions introduce computational complexity that provides no benefit when d=3, revealing fundamental algorithm-dimensionality mismatch.

\subsection{Precision Analysis}

\begin{table}[H]
    \footnotesize
    \centering
    \caption{Similar Images Found and Precision Analysis (Similarity threshold = 40.0)}
    \label{tab:precision_analysis}
    \setlength{\tabcolsep}{2pt}
    \begin{tabularx}{\columnwidth}{l >{\centering\arraybackslash}X >{\centering\arraybackslash}X >{\centering\arraybackslash}X >{\centering\arraybackslash}X >{\centering\arraybackslash}X}
        \toprule
        \textbf{Algorithm} & \textbf{1K} & \textbf{5K} & \textbf{10K} & \textbf{25K} & \textbf{26K} \\
        \midrule
        Linear Search & 303 & 1483 & 3109 & 7992 & 8380 \\
        Hash Search & 387 & 1907 & 3987 & 10203 & 10697 \\
        Hash Dynamic & 303 & 1483 & 3109 & 7992 & 8380 \\
        Quadtree & 303 & 1483 & 3109 & 7992 & 8380 \\
        Octree & 303 & 1483 & 3109 & 7992 & 8380 \\
        LSH & 288 & 1423 & 2997 & 7668 & 8039 \\
        M-Tree & 282 & 1265 & 2694 & 7272 & 7650 \\
        \bottomrule
    \end{tabularx}
    \\[0.1cm]
    \footnotesize\textit{Values represent number of similar images found within threshold=40.0. Linear Search serves as ground truth baseline for exact similarity matching.}
\end{table}
\vspace{-0.3cm}

\textbf{Precision Analysis Findings:}

\textbf{Exact Matching Cluster:} Four classical structures (Linear Search, Hash Dynamic, Quadtree, Octree) achieve identical 100\% precision with 8380 results for the 26K dataset. This validates correct threshold application and establishes classical spatial indexing as ground truth for RGB similarity search. The exact match across diverse algorithmic approaches confirms implementation correctness and demonstrates that simple geometric partitioning suffices for precise RGB similarity.

\textbf{Hash Search Systematic Expansion:} Consistently finds 1.28× more images (10697 vs 8380) due to dynamic cell sizing (cell=30.0) capturing neighboring color regions beyond exact threshold boundaries. This systematic 27.7\% over-inclusion persists across all scales, suggesting intentional broader matching rather than algorithmic error—potentially valuable for exploratory similarity applications accepting relaxed precision.

\textbf{LSH Approximation Trade-off:} Achieves 95.9\% recall (8039/8380 results), demonstrating the inherent precision-performance trade-off of probabilistic hashing. The 4.1\% precision loss stems from random projection quantization and probabilistic bucket collisions. However, this precision sacrifice provides no performance benefit—LSH search time (4.62ms) significantly exceeds exact classical methods (Octree: 2.12ms), invalidating LSH's fundamental value proposition of trading precision for speed in RGB 3D space.

\textbf{M-Tree Precision Anomaly:} Achieves only 91.3\% recall (7650/8380 results), contradicting M-Tree's theoretical guarantee of exact metric search. The 8.7\% precision loss likely stems from aggressive covering radius pruning that eliminates valid candidates near threshold boundaries. This precision degradation provides no compensating performance advantage (2.46ms vs 2.12ms Octree), rendering M-Tree unsuitable for applications requiring threshold compliance. The precision loss suggests implementation challenges in balancing covering radius optimization with search completeness.


\section{Discussion}
\label{sec:discussion}

The experimental results comparing seven data structures reveal critical insights into the applicability of high-dimensional indexing methods to low-dimensional RGB similarity search, challenging assumptions about algorithmic sophistication universally translating to practical performance gains.

\subsection{Algorithm Selection Guidelines}

\textbf{Production Systems (>10K images):} Deploy Octree for optimal performance-precision balance (2.12ms search, 100\% precision, 5.48ms insertion). Quadtree provides competitive alternative (2.25ms search, 6.92ms insertion). Hash Dynamic offers slightly faster search (2.15ms) with identical precision.

\textbf{Moderate Datasets (1K-10K images):} Linear Search (1.00ms insertion, 2.32ms search, 100\% precision) provides excellent performance with minimal implementation complexity. The algorithmic simplicity and guaranteed precision outweigh marginal performance gains from sophisticated structures at this scale.

\textbf{Avoid for RGB Applications:} LSH (53.91ms insertion, 4.62ms search, 95.9\% recall) and M-Tree (11.85ms insertion, 2.46ms search, 91.3\% recall) provide no practical benefit over classical structures while sacrificing both performance and precision. These structures designed for high-dimensional spaces impose prohibitive overhead in RGB 3D.

\subsection{Scalability and Theoretical Divergence}

Classical structures demonstrate excellent scalability from 1K to 26K images. Octree insertion scales from 0.15ms to 5.48ms, with search scaling from 0.06ms to 2.12ms. LSH exhibits catastrophic failure: insertion scales from 2.34ms to 53.91ms, exceeding all classical structures. M-Tree insertion scales from 0.30ms to 11.85ms, providing no search advantage (2.46ms vs 2.12ms Octree) while losing 8.7\% precision.

LSH theory promises $O(n^\rho)$ sublinear query time for $\rho < 1$ in high dimensions, yet 26K search (4.62ms) is 2.18× slower than Octree (2.12ms). The fundamental mismatch: LSH design assumes d≥50 dimensions where random projections preserve distances; in RGB d=3, projection overhead dominates without benefit. M-Tree's 91.3\% recall contradicts its exact search guarantee, revealing challenges in covering radius computation when threshold=40.0 spans significant color space. Octree's geometric pruning aligns naturally with RGB cube geometry, enabling effective subtree elimination.


\section{Conclusion}
\label{sec:conclusion}

This comprehensive empirical study challenges conventional wisdom about data structure selection for RGB image similarity search by comparing classical spatial indexing methods against modern high-dimensional structures. Our evaluation of seven distinct approaches across 1K to 26,179 real images from the Animals-10 Kaggle dataset reveals that specialized high-dimensional structures provide no benefit in RGB 3D space while imposing significant performance and precision penalties.

\subsection{Key Research Contributions}

\textbf{Classical Structure Dominance:} Octree achieves optimal performance, combining fastest search time (2.12ms for 26K images) with 100\% precision (8380 exact results). Quadtree (2.25ms search) and Hash Dynamic (2.15ms search) provide competitive alternatives with identical precision. Simple geometric partitioning demonstrates that algorithmic sophistication does not guarantee practical superiority.

\textbf{LSH Catastrophic Failure:} Locality-Sensitive Hashing exhibits complete failure in RGB 3D: 53.91ms insertion (53.9× slower than Linear), 4.62ms search (2.18× slower than Octree), and 95.9\% recall (4.1\% precision loss). The algorithm-dimensionality mismatch—LSH optimizes for d≥50 dimensions where random projections preserve distances; in RGB d=3, projection overhead dominates without benefit—reveals critical limitations in applying high-dimensional methods to low-dimensional problems.

\textbf{M-Tree Precision Sacrifice:} M-Tree's triangle inequality pruning, theoretically guaranteeing exact metric search, paradoxically achieves only 91.3\% recall (8.7\% precision loss, 7650/8380 results) while providing no search advantage (2.46ms vs 2.12ms Octree). The 11.85ms insertion overhead and precision degradation render M-Tree unsuitable for RGB similarity applications requiring threshold compliance. This precision loss contradicts M-Tree's theoretical guarantees and suggests implementation challenges in balancing covering radius optimization with search completeness.

\textbf{Bidirectional Curse of Dimensionality:} Results demonstrate that high-dimensional indexing structures suffer bidirectionally—degrading not only in truly high-dimensional spaces (their original design motivation) but also imposing unnecessary complexity in low-dimensional spaces. RGB 3D represents a "sweet spot" where simple spatial partitioning (Octree, Quadtree) dominates both performance and precision. Structures optimized for curse-of-dimensionality mitigation ironically become victims of reverse curse effects when applied outside their design context.

\textbf{Linear Search Competitive Reality:} Simple sequential scanning (1.00ms insertion, 2.32ms search, 100\% precision) remains competitive at 26K scale, demonstrating that algorithmic sophistication does not guarantee practical superiority. Cache-friendly sequential access patterns often outperform pointer-chasing tree traversal for moderate datasets, emphasizing that implementation details and hardware characteristics significantly influence practical performance rankings.

\subsection{Practical Implications and Future Work}

\textbf{Recommendations:} Deploy Octree for RGB similarity systems with >10K images (optimal performance-precision balance). Use Linear Search for <1K datasets (implementation simplicity, guaranteed precision). Avoid LSH and M-Tree for RGB applications—their theoretical high-dimensional advantages translate to practical disadvantages in RGB 3D space.

\textbf{Future Directions:} (i) Systematic investigation to identify the dimensionality threshold where high-dimensional structures begin outperforming classical spatial indexing (empirical studies spanning d=3 RGB through d=128 SIFT to d=512 deep embeddings); (ii) LSH parameter optimization for low dimensions (investigating whether reduced L and k values decrease overhead while maintaining search quality); (iii) M-Tree precision investigation (analyzing covering radius computation and node splitting to restore exact search guarantees); (iv) Alternative color spaces (evaluating LAB, LUV for different spatial locality characteristics); (v) GPU acceleration potential (determining whether massive parallelism can overcome LSH's sequential overhead).


\section*{Author Contributions}
\textbf{Luan Carrieiros:} Project lead, Hash Dynamic, LSH implementation, OpenCV integration. \textbf{Diego Rocha:} Octree/Quadtree. \textbf{Iago Ribeiro:} Linear Search, theoretical analysis. \textbf{Bernardo Temponi:} Hash Search. \textbf{Arthur Moraes:} M-Tree, optimization, CREATE→TEST→DESTROY pattern.

\section*{Code Availability}
Complete source code available at: \url{https://github.com/LuanCarrieiros/PAA}

% BIBLIOGRAFIA COM NATBIB (já incluído no sbc2023.cls)
\bibliographystyle{sbc}
\bibliography{refs}

% REFERENCIAS MANUAIS (para garantir que apareçam)
\begin{thebibliography}{99}

\bibitem{indyk1998approximate}
Indyk, P. and Motwani, R. (1998).
Approximate nearest neighbors: Towards removing the curse of dimensionality.
In \textit{Proceedings of the Thirtieth Annual ACM Symposium on Theory of Computing (STOC '98)}, pages 604--613. ACM, New York, NY, USA.

\bibitem{ciaccia1997m}
Ciaccia, P., Patella, M., and Zezula, P. (1997).
M-tree: An efficient access method for similarity search in metric spaces.
In \textit{Proceedings of the 23rd International Conference on Very Large Data Bases (VLDB '97)}, pages 426--435. Morgan Kaufmann Publishers Inc., San Francisco, CA, USA.

\bibitem{kaggle_animals10}
Corrado, A. (2018).
Animals-10.
Dataset available at: \url{https://www.kaggle.com/datasets/alessiocorrado99/animals10}. Accessed: 2025-11-21.

\bibitem{opencv_library}
Bradski, G. (2000).
The OpenCV library.
\textit{Dr. Dobb's Journal of Software Tools}, 120:122--125.

\bibitem{8394553}
Lee, C.-H., Lu, H.-y., and Horng, J.-H. (2018).
Color quantization by hierarchical octa-partition in RGB color space.
In \textit{2018 IEEE International Conference on Applied System Invention (ICASI)}, pages 147--150.

\bibitem{9035427}
Saad, A.-M.~H.~Y., Abdullah, M.~Z., Alduais, N.~A.~M.~M., and Sa'ad, H.~H.~Y. (2020).
Impact of spatial dynamic search with matching threshold strategy on fractal image compression algorithm performance: Study.
\textit{IEEE Access}, 8:52687--52699.

\end{thebibliography}

\end{document}